\documentclass[11pt]{article}

\usepackage{amsmath, amssymb, amsthm}
\usepackage[margin=1in]{geometry}

\title{Norms: Definition and Examples}
\author{}
\date{}

\begin{document}

\maketitle

\section{Definition: Norm}

Let \(V\) be a vector space over the field of real numbers
\(\mathbb{R}\) (or complex numbers \(\mathbb{C}\)). A \textbf{norm} on
\(V\) is a function
\[
  \|\cdot\| : V \to \mathbb{R}
\]
that assigns a strictly non-negative real number to every vector in
\(V\), satisfying the following four axioms for all vectors
\(x, y \in V\) and all scalars \(c \in \mathbb{R}\):

\begin{enumerate}
  \item \textbf{Positive Definiteness:} \(\|x\| \ge 0\), and
        \(\|x\| = 0\) if and only if \(x = 0\).
        \emph{(The length of a vector is positive, and only the zero
        vector has zero length.)}
  \item \textbf{Absolute Homogeneity (Scaling):}
        \(\|c x\| = |c| \, \|x\|\).
        \emph{(Scaling a vector by \(c\) scales its length by the
        absolute value of \(c\).)}
  \item \textbf{Subadditivity (Triangle Inequality):}
        \(\|x + y\| \le \|x\| + \|y\|\).
        \emph{(The length of the sum of two vectors is at most the sum
        of their lengths.)}
  \item \textbf{Non-negativity of the norm:} \(\|x\| \ge 0\)
        \emph{(This is technically implied by 1 and 2, but often listed
        explicitly).}
\end{enumerate}

A vector space \(V\) equipped with a norm is called a \textbf{normed
vector space}, denoted \((V, \|\cdot\|)\).

\hrulefill

\section{The \(L^p\) Norms (The Family of Norms)}

In finite-dimensional spaces like \(\mathbb{R}^n\), the most common
norms are the \(L^p\) norms. For a vector
\(x = (x_1, x_2, \dots, x_n) \in \mathbb{R}^n\) and a real number
\(p \ge 1\), the \(L^p\) norm is defined as:
\[
  \|x\|_p = \left( \sum_{i=1}^n |x_i|^p \right)^{1/p}
\]

\subsection{Example 1: The \(L^\infty\) Norm (Maximum Norm)}

The \(L^\infty\) norm is the limit of the \(L^p\) norm as
\(p \to \infty\). It is defined as the maximum absolute value of the
components of the vector:
\[
  \|x\|_\infty = \max_{1 \le i \le n} |x_i|
\]

\begin{itemize}
  \item \textbf{Intuition:} This measures the ``worst-case'' or largest
        single component of a vector.
  \item \textbf{In Discrepancy Theory:} If you have a
        set system \(\mathcal{S} = \{S_1, \dots, S_m\}\) and a coloring
        \(\chi\), the discrepancy vector is
        \(v = (\chi(S_1), \dots, \chi(S_m))\).
        The \(L^\infty\) norm is \(\max_i |\chi(S_i)|\). Minimizing
        this norm over all colorings \(\chi\) gives the standard
        definition of discrepancy:
        \(\min_\chi \max_i |\chi(S_i)|\).
\end{itemize}

\subsection{Example 2: The \(L^2\) Norm (Euclidean Norm)}

The \(L^2\) norm corresponds to \(p=2\). It is the standard
``straight-line'' distance in Euclidean space:
\[
  \|x\|_2 = \sqrt{ \sum_{i=1}^n |x_i|^2 }
          = \sqrt{x_1^2 + x_2^2 + \dots + x_n^2}
\]

\begin{itemize}
  \item \textbf{Intuition:} This measures the standard length of a
        vector.
  \item \textbf{In Discrepancy Theory:} The formula in your image is exactly the \(L^2\) norm applied to the
        discrepancy vector.
        \[
          D_2(\mathcal{S})
          = \min_\chi \sqrt{\chi(S_1)^2 + \dots + \chi(S_m)^2}
          = \min_\chi \bigl\| (\chi(S_1), \dots, \chi(S_m)) \bigr\|_2
        \]
        As your text notes, the \(L^2\) norm is ``friendlier to
        manipulate algebraically'' because it is a smooth function
        (unlike the max function), allowing mathematicians to use
        techniques like linear algebra, Fourier analysis, and calculus
        to find optimal colorings.
\end{itemize}

\subsection{Example 3: The \(L^1\) Norm (Taxicab / Manhattan Norm)}

The \(L^1\) norm corresponds to \(p=1\). It is the sum of the absolute
values of the components:
\[
  \|x\|_1 = \sum_{i=1}^n |x_i|
\]

\begin{itemize}
  \item \textbf{Intuition:} This is the distance a taxicab would drive
        in a grid-based city like Manhattan (moving only horizontally
        and vertically).
\end{itemize}

\hrulefill

\section{A Concrete Numerical Example}

Let \(v = (3, -4, 0, 1)\) be a vector in \(\mathbb{R}^4\).

\begin{itemize}
  \item \textbf{\(L^1\) Norm:}
        \(\|v\|_1 = |3| + |-4| + |0| + |1| = 3 + 4 + 0 + 1 = 8\)
  \item \textbf{\(L^2\) Norm:}
        \(\|v\|_2 = \sqrt{3^2 + (-4)^2 + 0^2 + 1^2}
                  = \sqrt{9 + 16 + 0 + 1}
                  = \sqrt{26} \approx 5.099\)
  \item \textbf{\(L^\infty\) Norm:}
        \(\|v\|_\infty = \max(|3|, |-4|, |0|, |1|) = 4\)
\end{itemize}

Notice that for any vector,
\(\|v\|_\infty \le \|v\|_2 \le \|v\|_1\). This relationship is a
fundamental result in normed spaces: all norms on a finite-dimensional
vector space are equivalent (they bound each other up to a constant
factor depending on the dimension).

\end{document}
